\section{Transparencias de distribución}

\begin{frame}{Spotify: acceso y ubicación}
  \small
  \begin{tabular}{@{}lp{0.70\textwidth}@{}}
    \textbf{Transparencia} & \textbf{Ejemplo y alcance} \\
    \hline
    Acceso &
      Desde la app, el usuario puede reproducir música
      tras activar la calidad sin pérdida. Spotify presenta la calidad
      como una preferencia de reproducción \cite{spotifyLossless2025}. \\[0.7em]
    Ubicación &
      La app accede al servicio sin que el usuario elija una región.
      El perímetro de Spotify recibe las solicitudes y distribuye el
      tráfico mediante regiones de nube \cite{spotifyIncident2025}. \\
  \end{tabular}
\end{frame}

\begin{frame}{Spotify: replicación y fallos}
  \small
  \begin{tabular}{@{}lp{0.70\textwidth}@{}}
    \textbf{Transparencia} & \textbf{Ejemplo y alcance} \\
    \hline
    Replicación &
      Replicación entre servidores: sin confirmar en las fuentes
      consultadas. \emph{Offline Backup} documenta el uso de pistas en
      la caché del dispositivo \cite{spotifyOffline2024}. \\[0.7em]
    Fallos &
      \emph{Offline Backup} permite reproducir pistas en caché cuando
      se pierde la conexión \cite{spotifyOffline2024}. La caída del
      perímetro del 16 de abril de 2025 interrumpió la mayor parte
      del tráfico \cite{spotifyIncident2025}. \\
  \end{tabular}

  \medskip
  {\footnotesize Contexto técnico sobre caché y distribución de medios
  \cite{viola2023media,van2023caching,nguyen2023edge}.}
\end{frame}
